%% kapa: exam=1 round=32 year=2021 period=1 form=A subject=civil_law subject_ko=민법 seq=33 number=33
%% source: Q-Net 감정평가사 기출문제 (cst00309) artlSeq=5208397 file=extracted/1008/5208397/A형/2021년도 제32회 감정평가사 1차 1교시 A형.pdf
%% conversion: kordoc 4.24.51 -> tools/gen1.py (auto); review=proofread-sonnet; flags=none
\begin{question}{33}
분묘기지권에 관한 설명으로 옳지 않은 것은? (다툼이 있으면 판례에 따름)
\begin{choices}
  \item 분묘기지권을 시효취득하는 경우에는 특약이 없는 한 지료를 지급할 필요가 없다.
  \item 「장사 등에 관한 법률」이 시행된 후 설치된 분묘에 대해서는 더 이상 시효취득이 인정되지 않는다.
  \item 분묘기지권의 시효취득을 인정하는 종전의 관습법은 법적 규범으로서의 효력을 상실하였다.
  \item 분묘기지권이 인정되는 분묘를 다른 곳에 이장하면 그 분묘기지권은 소멸한다.
  \item 분묘가 일시적으로 멸실되어도 유골이 존재하여 분묘의 원상회복이 가능하다면 분묘기지권은 존속한다.
\end{choices}
\end{question}
\answer{3}
